\chapter{How a big number is stored}
\label{ch:repr}

\section{Limbs: digits in base $2^{64}$}

On paper we write numbers in base 10: each digit is $0$--$9$ and each position is
worth ten times the one to its right. A computer can do the same thing with a much
bigger base. BigMath uses base $B = 2^{64}$: each ``digit'', called a \emph{limb}, is a
whole 64-bit machine word and can be anything from $0$ to
$18\,446\,744\,073\,709\,551\,615$.

\begin{keyidea}
A \code{BigInteger} is a list of 64-bit limbs plus a sign. The value is
$a_0 + a_1 B + a_2 B^2 + \dots$ with $B = 2^{64}$: exactly like decimal digits, but each
``digit'' is a 64-bit word.
\end{keyidea}

The limbs are stored \emph{little-endian}: the least significant limb comes first,
at index 0.

\begin{figure}[ht]
\centering
\begin{tikzpicture}
  \foreach \i/\v in {0/a_0,1/a_1,2/a_2,3/a_3} {
    \node[limb,minimum width=18mm] (L\i) at (\i*1.9,0) {$\v$};
    \node[lbl,below=1pt of L\i] {index \i};
  }
  \node[lbl,above=2pt of L0] {$\times B^0$};
  \node[lbl,above=2pt of L1] {$\times B^1$};
  \node[lbl,above=2pt of L2] {$\times B^2$};
  \node[lbl,above=2pt of L3] {$\times B^3$};
  \draw[arr] (-0.9,-0.95) -- node[below,lbl]{memory order: least significant first} (6.6,-0.95);
\end{tikzpicture}
\caption{A four-limb number in memory. Limb $i$ is worth $B^i = 2^{64 i}$.}
\label{fig:limbs}
\end{figure}

\begin{example}[: writing $2^{130} + 5$ in limbs]
Divide by $B = 2^{64}$ repeatedly, as you would to convert to any base:
\begin{itemize}
  \item $2^{130} + 5 = 2^{66}\cdot B + 5$, so $a_0 = 5$ and the rest is $2^{66}$.
  \item $2^{66} = 4 \cdot B + 0$, so $a_1 = 0$ and the rest is $4$.
  \item $4 < B$, so $a_2 = 4$.
\end{itemize}
The limbs are $[5,\ 0,\ 4]$. Check: $5 + 0\cdot 2^{64} + 4 \cdot 2^{128} = 2^{130} + 5$.
\end{example}

\subsection{Why base $2^{64}$?}

\begin{itemize}
  \item \textbf{The hardware already does it.} A 64-bit processor multiplies two
        64-bit words and returns the full 128-bit result in one or two
        instructions (\code{MUL}+\code{UMULH} on ARM, \code{MUL} on x86). So a
        limb-by-limb product is as cheap as it can be.
  \item \textbf{Fewer limbs means less work.} Many algorithms cost about $n^2$
        limb operations for $n$-limb inputs. Halving the number of limbs
        (base $2^{64}$ instead of $2^{32}$) quarters that work. The library used
        32-bit values in 64-bit slots for years; moving to full 64-bit limbs was
        the biggest single speedup in its history. The old layout is still
        available with \code{-DBIGMATH\_LIMB\_64=0} for comparison.
  \item \textbf{Binary operations are free.} Shifts, bit tests, and conversion
        to hexadecimal work directly on the stored bits.
\end{itemize}

The price is decimal conversion: turning limbs into decimal digits is real work
(\cref{ch:conv}). Storing numbers in base $10^{18}$ would make printing trivial but
every arithmetic operation slower; arithmetic happens far more often than printing.

\subsection{Why little-endian?}

Carries flow from the least significant limb to the most significant. With the low
limb at index 0, every loop runs forward through memory, which is what processors
prefetch best. A result that grows by one limb (a final carry) just gets one more
element appended at the end; nothing has to move.

\section{Sign and magnitude}

The sign is stored separately as a flag; the limbs hold the absolute value. So $-5$ is
the limbs $[5]$ with the sign flag set. Algorithms work on the magnitudes and the
operator layer (\file{include/biginteger/ops/}) works out the sign of the result:
for example, the product is negative when exactly one input is negative.

Division truncates toward zero, as in C and C++: the quotient is rounded toward zero
and the remainder has the sign of the dividend.

\begin{center}
\begin{tabular}{rrrr}
\toprule
$a$ & $b$ & $a / b$ & $a \% b$ \\
\midrule
$7$ & $2$ & $3$ & $1$ \\
$-7$ & $2$ & $-3$ & $-1$ \\
$7$ & $-2$ & $-3$ & $1$ \\
$-7$ & $-2$ & $3$ & $-1$ \\
\bottomrule
\end{tabular}
\end{center}

In every row $b \cdot (a/b) + (a \% b) = a$.

\section{Small numbers stay inside the object}
\label{sec:smallops}

Most programs that use big integers also use many \emph{small} ones: counters,
loop bounds, intermediate values that happen to be short. If every value needed a
heap allocation, copying the number $3$ would cost as much as calling \code{malloc}.

BigMath therefore keeps numbers of up to four limbs (256 bits,
\code{BIGMATH\_INLINE\_LIMBS}) \emph{inside} the \code{BigInteger} object. Only longer
numbers use a \code{std::vector}. The rule is strict: after trimming leading zero
limbs, a value is stored inline exactly when it has at most four limbs. So copying a
small value never allocates.

\begin{figure}[ht]
\centering
\begin{tikzpicture}
  \node[box,minimum width=60mm,minimum height=14mm,anchor=west] (obj) at (0,0) {};
  \node[lbl,anchor=south west] at (obj.north west) {\code{BigInteger} object};
  \foreach \i in {0,1,2,3} \node[limb,minimum width=11mm] at (0.75+\i*1.2,0) {$a_\i$};
  \node[lbl,anchor=west] at (obj.east) {\quad $\le 4$ limbs: stored inline, no heap};
  \node[box,minimum width=60mm,minimum height=14mm,anchor=west] (obj2) at (0,-2.2) {};
  \node[limb,minimum width=20mm] (p) at (1.3,-2.2) {pointer};
  \node[lbl] at (4.1,-2.2) {sign, length};
  \foreach \i in {0,...,7} \node[limb,minimum width=9mm] (h\i) at (7.5+\i*0.9,-3.4) {\scriptsize$a_{\i}$};
  \draw[arr] (p.south) |- (h0.west);
  \node[lbl,anchor=west] at (obj2.east) {\quad $\ge 5$ limbs: heap vector};
\end{tikzpicture}
\caption{Inline storage for short values, heap storage for long ones.}
\end{figure}

The operators go one step further. When both operands have at most 16 limbs
(\code{BIGMATH\_SMALL\_OPS\_MAX\_LIMBS}), \code{+}, \code{-}, \code{*}, \code{/}
and \code{\%} compute the result in a buffer on the stack using the limb kernels of
\cref{ch:basics}, then build the result object directly. No temporary vectors are
created; the only possible allocation is the result itself when it has more than four
limbs.

\begin{perfnote}
Before inline storage, a one-limb product took 5.7 times as long as GMP's, and copying
a small value 5.3 times as long. After it, the product takes 2.1 times GMP's time and a
copy 0.12 times (it is cheaper than GMP's because GMP has to allocate). For 1--16 limb
operands the cost is now dominated by the arithmetic, not by memory management.
\end{perfnote}

\begin{practice}
You do not need to do anything to benefit: the fast path is automatic. But it does
mean that code full of small \code{BigInteger} temporaries is fine; there is no reason
to fall back to \code{int64\_t} ``for speed'' unless you are in a very tight loop.
\end{practice}

\section{Two layers: operators and algorithms}

The code is split into two layers, and the split matters when you read the rest of
this guide.

\begin{figure}[ht]
\centering
\begin{tikzpicture}[node distance=7mm]
  \node[box,minimum width=95mm] (ops) {\textbf{Operators} \quad \file{ops/}\\ \footnotesize sign handling, small-operand stack path, \code{operator*}, \code{operator/}, \code{ToString}};
  \node[box,minimum width=95mm,below=of ops] (disp) {\textbf{Dispatchers} \quad \file{src/algorithms/*.cpp}\\ \footnotesize look at the sizes and pick an algorithm};
  \node[box,minimum width=95mm,below=of disp] (alg) {\textbf{Algorithms} \quad \file{algorithms/multiplication/}, \file{algorithms/division/}\\ \footnotesize schoolbook, Karatsuba, Toom-3, NTT, Knuth D, Burnikel--Ziegler, Newton};
  \node[box,minimum width=95mm,below=of alg] (ker) {\textbf{Limb kernels} \quad \file{common/LimbKernels.h}\\ \footnotesize add, subtract, multiply-by-one-limb, shift on raw arrays};
  \draw[arr] (ops) -- (disp);
  \draw[arr] (disp) -- (alg);
  \draw[arr] (alg) -- (ker);
  \draw[arr] (alg.east) -- ++(0.8,0) |- node[pos=0.25,right,lbl,align=left]{division and\\conversion call\\the dispatcher\\for sub-products} (disp.east);
\end{tikzpicture}
\caption{Layers of the library. The upward arrow matters: a division step or a decimal
conversion calls the multiplication dispatcher for its sub-products, so each sub-product
gets the best algorithm for its own size.}
\label{fig:layers}
\end{figure}

Arithmetic operators borrow read-only limb ranges through \code{BigInteger::Limbs()}.
A view holds an address and a length. It can cover inline limbs, heap-backed limbs,
or a slice of either while the original storage stays valid. Kernels keep their
writable output and workspace separate. The upward arrow in \cref{fig:layers} lets
the dispatcher choose the method for each sub-product.
